\section{Verification and the Agent Harness}
\label{sec:harness}

\subsection{Exact verifiers}
\label{sec:verifiers}
Four verifiers are exact on their domains. \emph{Mathematics}: numeric
equality with a reference and symbolic equality through SymPy. \emph{Code}:
hidden unit tests run in the sandbox, with per-assert credit protected by a
per-run nonce so that a failing candidate still gets a graded score.
\emph{Structured output}: JSON Schema validation. \emph{Proofs}: Lean~4
\cite{moura2021lean4} under the frozen-statement protocol of Open Math Harness
\cite{openmathharness2026}. The statement is frozen on the service side; the
candidate is refused on text alone for forbidden tokens (\texttt{sorry},
\texttt{axiom}, \texttt{native\_decide}, commands that run code during
elaboration) or for not reproducing the statement byte for byte; Lean then
checks a file assembled from the frozen preamble, the candidate's body, an
unguessable marker, the axioms report and the fully elaborated type, which must
equal the pristine statement's character for character. Against real Lean
4.35.0 the protocol accepts three correct proofs and rejects every known
attack at the layer assigned to it (Table~\ref{tab:reallean}); the decisive
case is a candidate that shadows the addition instance and closes the
statement with \texttt{rfl}, which only the comparison of elaborated types
rejects. A check takes about 0.6\,s. Requests are refused when no toolchain is
configured, so a Lean task can never be labelled verified by a fallback.

\begin{table}[t]
\centering
\caption{The Lean protocol against real Lean 4.35.0, for the statement \texttt{a + b = b + a}.}
\label{tab:reallean}
\scriptsize
\begin{tabular}{@{}lll@{}}
\toprule
Candidate & Verdict & Decided by \\
\midrule
\texttt{by omega} & verified & Lean \\
\texttt{Nat.add\_comm a b} & verified & Lean \\
helper lemma, then the theorem & verified & Lean \\
\texttt{by sorry} & partial (0.5) & Lean elaborates; \texttt{sorry} \\
\texttt{by rfl} & errors (0.05) & Lean error \\
\texttt{sorryAx \_ false} & rejected & axioms report \\
\texttt{axiom cheat : False} & rejected & textual rule \\
shadowed \texttt{HAdd} instance, \texttt{rfl} & rejected & elaborated type \\
\bottomrule
\end{tabular}
\end{table}

The soft verifiers are the rubric judge and sample consensus; both yield an
assurance level below \texttt{verified}. Measured rates for all verifiers are
in Section~\ref{sec:eval:gate}.

\subsection{Best-first repair}
Every verified attempt of the deliberate path becomes a node of a tree, scored
from the verifier's verdict alone (1 solved; for a partial result with $C$
credited and $L$ open steps, $0.5+0.4\,C/(C+L)$; 0.05 for errors; 0 for
invalid). The next attempt repairs the node a policy chooses and sees that
node's evidence and, as data, what earlier repairs of it did. With one
candidate explored at a time the policy is: highest score, ties to the fewest
failures, then to the latest node; a node stays continuable until it has $F$
children (default 2), which lets the search return to a parent when a repair
made things worse.

\begin{proposition}[Conservative extension] \textnormal{[exact]}
\label{prop:h1}
If every verdict is all-or-nothing (every failed candidate scores 0.05 and has
one failure), the best-first policy repairs the latest candidate at every round,
exactly as a plain retry loop does, for any fanout $F\ge1$ along a chain of
repairs.
\end{proposition}
\begin{proof}
Along a chain each new node is a child of the previous one and has no children
of its own, so it is eligible. All failed nodes tie on score and failures; the
tie goes to the latest node. The root scores 0 and is chosen only when no
eligible node scores more, which a valid failed candidate (0.05) prevents. \qed
\end{proof}
The policy therefore changes behaviour only when per-assert credit, Lean's
graded verdicts or an invalid patch distinguish candidates.

\subsection{The code agent}
A bounded agent works on a workspace of flat files. Each reply holds one
action: a patch in Codex's \texttt{apply\_patch} language
\cite{openaicodex2026}, a whole file, or a \texttt{run} of a workspace file in
the sandbox. The exact verifier runs after every change and a passing
workspace ends the run at once; the model's own claim of completion never
does. The best-scoring workspace is returned when nothing passes. Model calls
are bounded by a step limit (default 6, at most 12) and the run by a
wall-clock cap that holds inside a model call, as Lemma~\ref{lem:termination}
requires of every loop.

\textbf{Where a step spends its time.} To first order
\begin{equation}
T_{\text{step}} = \frac{n_{\text{prompt}}-n_{\text{cached}}}{r_p} + \frac{n_{\text{dec}}}{r_d} + T_{\text{tools}},
\label{eq:stepcost}
\end{equation}
with prefill rate $r_p$ and decode rate $r_d$. \Stwo\ prefills about 31
tokens/s and decodes 6 to 8: a cold 1{,}969-token prompt took 64\,s to
prefill, and the same prompt with one more exchange appended took 3\,s.
Whether the cache is hit decides the cost of a step.

\textbf{An append-only transcript.}
\begin{proposition}[Prefix reuse]\label{prop:prefix}
Let a server reuse the longest common prefix between a request and the token
sequence cached for the previous request of the run. If the files are part of
the first message and are rewritten each step, then after a step that changed
a file the reusable prefix ends at the first changed byte, and a run of $K$
steps prefills $\Theta(K\cdot L)$ tokens for prompts of length $L$. If the
transcript only grows, step $k$ prefills only the previous reply and its
observation, and a run prefills each transcript token once: $\Theta(L_K)$ for
a final transcript of length $L_K$, plus one task message per compaction.
\end{proposition}
The transcript is therefore never rewritten: a change of state is a new
observation carrying the edited files or the verifier's evidence. Its size is
metered from the server's reported prompt tokens plus an estimate of what was
appended since. When it exceeds its budget, older observations are first cut
to a head and a tail, and only if that is not enough is the transcript
compacted into a \emph{harness-written} state (the task, the current files,
the last verdict and an ``already tried'' list). A model-written summary, as
Codex uses, would cost minutes at these decode rates, and the facts it needs
are already structured.

\textbf{Actions under a grammar.} On a backend that accepts a decoding
grammar, each step's grammar admits only that step's legal actions: update and
delete targets are the workspace's files, a hunk header cannot start with a
line range, a \texttt{run} can name only an existing file, and the last step
offers only edits. A reasoning model needs one more rule. A grammar applies
from the first sampled token, and with a plain grammar Qwen3-4B-Thinking was
forced to emit the answer in place of its reasoning; the grammar therefore
opens with a free reasoning span up to the closing tag:
\begin{lstlisting}
root  ::= think action
think ::= ( [^<] | "<" [^/] | ... )* "</think>" [ \n]*
action ::= patch | whole | run
whole ::= "```python\n" codeline* "```"
\end{lstlisting}
The reasoning stays unconstrained \cite{tam2024speakfreely}. On a backend
without grammars the harness reads free text with a lenient parser (a patch
may start at its first file section, and unified-diff line ranges in a hunk
header are dropped).

\textbf{Which edits to offer.} Let a whole file cost $W$ answer tokens, a
patch $P$, and let a patch fail to apply with probability $f$. The expected
answer tokens of a repair are about $P+fW$, so a patch pays when
\begin{equation}
f<1-P/W .
\label{eq:patch}
\end{equation}
Small models fail patch formats often, and a patch that applies can still
break the file. The harness therefore offers the whole file while it is
missing or at most 80 lines long, and a patch or the whole file above that.

\textbf{Guards.} (i)~After an edit every changed Python file is compiled in
the sandbox, never executed; a file that does not compile is not applied, and
the model sees the error with the lines around it (37\,ms per edit).
(ii)~A change judged below the best version so far is rolled back.
(iii)~A failed step's signature is the hidden tests' evidence, or the action
and its first line of feedback; the second consecutive identical failure
carries a nudge, the fourth ends the run as \texttt{stuck}. The same action
text sent a third time draws a notice before it fails. After a step that
reproduced an already-judged workspace the transcript is compacted and the
next step is sampled hotter. (iv)~An edit in a reply that ended on the token
budget is refused. (v)~A failed model call is retried twice with backoff.
(vi)~The first step of a \Stwo\ run uses a 256-token reasoning budget and
later steps 512.

\textbf{A cascade.} The gate's logic applies to the agent as well. The first
step goes to \Sone\ as one request with four choices (the prompt prefilled
once, the choices decoded in one batch at temperature 0.7); each candidate is
applied to its own copy of the files and judged by the exact verifier; the
first that passes ends the run. When none passes, the best-scoring candidate
becomes the state and \Stwo\ continues from step two.
